\chapter{Wstęp}
\label{chap:wstep}

Celem niniejszego projektu było zaprojektowanie i zaimplementowanie zaawansowanego systemu biometrycznego czasu rzeczywistego, który integruje dwie kluczowe technologie: rozpoznawanie twarzy oraz analizę emocji. System został opracowany w języku Python z wykorzystaniem bibliotek open-source, jednak kluczowe komponenty, takie jak logika biznesowa, zarządzanie danymi oraz algorytmy interpretacji cech biometrycznych, stanowią autorską implementację.

Projekt wpisuje się w obszar zagadnieień związanych z bezpieczeństwem systemów informatycznych oraz interakcją człowiek-komputer (HCI). Główne założenia funkcjonalne systemu obejmują:
\begin{itemize}
    \item Uwierzytelnianie użytkowników na podstawie unikalnych cech ich twarzy.
    \item Analizę mikroekspresji w celu określenia konkretnej emocji użytkownika w czasie rzeczywistym.
    \item Stworzenie spersonalizowanego modelu biometrycznego dla każdego użytkownika poprzez mechanizm kalibracji wykrywania emocji.
    \item Zapewnienie wysokiej wydajności działania dzięki zastosowaniu technik optymalizacyjnych.
\end{itemize}
